% TOP_PROBLEM: {"id": 3044, "title": "Wide partition conjecture", "category_id": 2, "category": {"id": 2, "name": "combinatorics", "display_name": "Combinatorics", "description": "Counting problems, graph theory, discrete structures.", "slug": "combinatorics", "order_index": 2, "created_at": "2026-07-31T15:26:25.671Z"}, "status": "open", "problem_number": "OPG-2063", "set_id": 12}
% FIRST_POSTED: 2026-10-01
% ATTEMPT_STATUS: unresolved
% UnsolvedMath ID 3044; source catalogue code OPG-2063
% !TeX program = lualatex
% Research attempt by GPT-6 Astra Ultra; no claim of a new complete solution.
\documentclass[11pt,a4paper]{article}
\usepackage[margin=25mm]{geometry}
\usepackage{fontspec}
\setmainfont{lmroman10-regular.otf}[BoldFont=lmroman10-bold.otf,
 ItalicFont=lmroman10-italic.otf,BoldItalicFont=lmroman10-bolditalic.otf]
\usepackage{amsmath,amssymb,mathtools}
\usepackage{unicode-math}
\setmathfont{latinmodern-math.otf}
\AtBeginDocument{\let\setminus\smallsetminus}
\usepackage{xurl,hyperref}
\hypersetup{colorlinks=true,urlcolor=blue}
\setcounter{secnumdepth}{0}
\setlength{\parindent}{0pt}
\setlength{\parskip}{5pt}
\setlength{\emergencystretch}{3em}
\begin{document}
\section*{Wide partition conjecture}\label{3044}
{\small UnsolvedMath ID 3044 · OPG-2063 · Combinatorics · OpenGarden}
\subsection{Definitions and mathematical statement}
An integer partition is a finite nonincreasing sequence
$\lambda=(\lambda_1,\ldots,\lambda_s)$ of positive integers.
A subpartition $\mu$ is obtained by selecting any submultiset of its parts
and putting the selected parts in nonincreasing order. Thus entire rows
are removed; individual parts are not shortened. Define the conjugate
partition by $\mu'_j=|\{i:\mu_i\ge j\}|$.

For partitions of the same integer, the dominance relation $\alpha\succeq\beta$
means $\sum_{i=1}^t\alpha_i\ge\sum_{i=1}^t\beta_i$ for every $t\ge1$,
with missing parts treated as zero. Say that $\lambda$ is wide if
$\mu\succeq\mu'$ for every subpartition $\mu$, including the empty one.

The diagram of $\lambda$ has cells $(i,j)$ with $1\le i\le s$ and
$1\le j\le\lambda_i$, with all rows left justified. Call $\lambda$ Latin
if its cells admit integer entries $T_{ij}$ such that row $i$ is a
permutation of $1,\ldots,\lambda_i$, and entries in each column are
pairwise distinct. No condition requires distinctness across different
columns beyond the row condition.

The wide partition conjecture states that for every partition $\lambda$,
\[
 \lambda\text{ is wide}\quad\Longleftrightarrow\quad
 \lambda\text{ is Latin}.
\]
The dominance tests are imposed on every subpartition, not only on
$\lambda$ itself. The constructive direction asks for one simultaneous
filling of the entire diagram from those inequalities; separate fillings
of individual rows are not sufficient.
\subsection{Short English statement}
Are exactly the partitions passing every subpartition dominance test those whose rows can be filled as permutations without repeated symbols in a column?
\subsection{Research attempt}
\paragraph{Scope and process.}
This is a GPT-6 Astra Ultra research attempt dated 1 October 2026, using
primary-source browsing and elementary mathematical checks. The canonical
definition above is retained. Results below are restricted results or
reductions, not a claim of a new solution to the entire catalogue question.
No formal proof assistant or external mathematical referee checked this text.


\paragraph{Literature and plan.}
Chow, Fan, Goemans and Vondr\'ak [R1] formulate the wide-partition
problem and prove partial results. Aharoni, Berger, Guo and Kotlar
[R2] prove a dual version; this is not automatically the desired
simultaneous filling theorem. We prove the necessary dominance
conditions directly and construct fillings for rectangles and for
every admissible partition with at most two rows.

\paragraph{Latin implies wide: a capacity proof.}
Let $\mu$ be any selected subpartition of a Latin partition, keeping
the selected rows in nonincreasing length order. The inherited filling
is still Latin. For any $t\ge1$, the first $t$ columns contain
\[
 \sum_i\min(t,\mu_i)=\sum_{j=1}^t\mu'_j
\]
cells. Symbol $s$ occurs exactly $\mu'_s$ times in the entire filling,
because it occurs once in each row of length at least $s$. Among the
first $t$ columns it can occur at most $t$ times, by column
distinctness. Hence the number of cells there is at most
\[
 \sum_{s\ge1}\min(t,\mu'_s)=\sum_{i=1}^t\mu_i.
\]
The last identity counts the boxes of the first $t$ rows by columns.
Comparing the two displays proves $\mu\succeq\mu'$ for every $t$.
Since the selected rows were arbitrary, the original partition is wide.
This proves one full direction of the conjecture without any unproved
matching or integrality assumption.

\paragraph{Rectangles.}
For $\lambda=(a^b)$, the first dominance inequality requires
$a\ge b$. If this holds, fill
\[
 T_{ij}=1+((i+j-2)\bmod a),\qquad 1\le i\le b,\quad1\le j\le a.
\]
Each row is a cyclic permutation of $[a]$. Within a column, distinct
row indices have distinct residues modulo $a$, since there are at
most $a$ rows. Thus the rectangle is Latin, and the preceding
necessity proof shows it is wide. If $a<b$ it is neither, because
the first column cannot have $b$ distinct entries drawn from $[a]$.
This settles the equivalence for every rectangular shape.

\paragraph{All shapes with two rows.}
Write $\lambda=(a,b)$ with $a\ge b\ge1$. The shape $(1,1)$ is
neither wide nor Latin. For $a\ge b\ge2$, put $1,2,\ldots,a$ in
the first row and $2,3,\ldots,b,1$ in the second. The lower row is
a derangement of $[b]$ relative to the first $b$ entries, so no
column repeats. For $b=1$ and $a\ge2$, put
$2,1,3,\ldots,a$ in the first row and $1$ in the second. This
also works. Hence every other two-row shape is Latin and therefore
wide. A one-row shape is filled in its natural order. These checks
include the row-length-one boundary that a single cyclic-shift formula
would otherwise miss.

\paragraph{Why the capacity argument does not construct a general filling.}
For each symbol $s$, one needs to place it once in every row of
length at least $s$ and at most once in every column. Such placements
for different $s$ must occupy disjoint cells. Solving these matching
problems separately does not enforce that last requirement. In the
square shape $(2,2)$, for example, each symbol separately can be
placed in the two diagonal cells, giving two individually valid
matchings which conflict at both cells. The shape is Latin, but this
demonstrates the logical failure of choosing arbitrary separate
matchings and superimposing them. The dominance inequalities above
are necessary capacities; an additional coordinated selection theorem
is needed to make them sufficient in all shapes.

\paragraph{Verification and final status.}
Subpartitions delete whole rows, exactly as required; no shortened
rows or arbitrary skew diagrams are introduced. The capacity count
uses the prescribed row alphabet $[\mu_i]$, not just rowwise
distinctness with arbitrary symbols. \textbf{UNRESOLVED}: necessity
and the stated infinite families are proved, but sufficiency for an
arbitrary wide partition remains the first missing step.

\subsection{Sources}
\begin{itemize}
\item[S1] ULAM.AI and the UnsolvedMath Contributors, supplied problems.json, record 3044 (OPG-2063), OpenGarden collection. Catalogue identity and source transcription; CC BY 4.0.
\url{https://huggingface.co/datasets/ulamai/UnsolvedMath}.
\item[S2] Open Problem Garden, Wide partition conjecture. Original problem and discussion; the mathematical formulation below is expanded from the supplied source record.
\url{http://www.openproblemgarden.org/op/wide_partition_conjecture}.

\item[R1] Timothy Y. Chow, C. Kenneth Fan, Michel X. Goemans and
Jan Vondr\'ak, \emph{Wide partitions, Latin tableaux, and Rota's basis
conjecture}, Advances in Applied Mathematics 31 (2003), 334--358.
\url{https://math.mit.edu/~goemans/PAPERS/wide.pdf}.
\item[R2] Ron Aharoni, Eli Berger, He Guo and Daniel Kotlar,
\emph{2-covers of wide Young diagrams}, arXiv:2311.17670v3,
revised 21 August 2025. Primary abstract checked for the scope of its
dual result.
\url{https://arxiv.org/abs/2311.17670}.

\end{itemize}
\end{document}
